\chapter{Merging ordered insertions}
\label{chap:merging-ordered-insertions}

Bringing two points together asks two different questions.
A Laurent expansion retains their separation as a small parameter.
A collision sets that separation to zero and must specify the state
left at the merged point.
The expansion at the end of
Chapter~\ref{chap:ordered-coefficient-bar} answers the first question.
We now construct maps answering the second for ordered coefficients
and a free algebra of composite states.

Order matters at the moment of collision.
The two branches of the binary coefficient algebra give two
restrictions to the same point.
With three labels, successive restrictions must retain the same
total order.
We will write the maps, the complexes recording the equations of
collision, and the changes of variables that prove this compatibility.
Their bar duals require an additional change of coefficients.
A binary example will show why that change cannot be omitted.

\section{Which charts survive a collision?}

Retain the polynomial ring $R_I=\mathbb C[z_i\mid i\in I]$ and
the coefficient algebra $C_I$ of
Chapter~\ref{chap:ordered-coefficient-bar}.
Its components use a nonempty collection of total-order charts.
A difference $z_i-z_{i'}$ is inverted exactly when that pair
changes order among those charts.
Compatible polynomial forms on their simplices define $C_I$.
The differential acts on the simplex variables and fixes all $z_i$.

Let $p:I\to J$ be a surjection of nonempty finite sets.
The labels in $p^{-1}(j)$ are to become one label $j$.
Choose an order $\tau_j$ on each such fibre.
An order $\sigma$ of $J$ expands to an order $e_p(\sigma)$ of $I$:
list the blocks according to $\sigma$, and list their members
according to $\tau_j$.
This is an injection between the two finite sets of orders.

Write $Z_j$ for the merged coordinates.
The coordinate substitution is
\[
 R_I\longrightarrow R_J,\qquad z_i\longmapsto Z_{p(i)}.
\]
It cannot be applied to every localization in $C_I$.
For example, an inverse of $z_1-z_2$ cannot be specialized when
$p(1)=p(2)$.
The selected charts remove exactly this difficulty.

For a nonempty collection $L$ of orders of $J$, consider the
component indexed by $e_p(L)$.
The relative order of two labels in one block is fixed on this
component. Their difference is therefore never inverted.
For labels in different blocks, their order varies precisely when
the order of their two blocks varies in $L$.
Consequently coordinate substitution defines
\[
 R_{I,e_p(L)}\longrightarrow R_{J,L}.
\]
Restrict a coefficient family to these selected components,
replace $t_{e_p(\sigma)}$ by $t_\sigma$, and make this substitution.
Face restrictions commute with both steps.

\begin{proposition}\label{prop:moi-coefficient-map}
The prescribed formulas define a unital morphism of commutative
differential graded algebras
\[
 k_p:C_I\longrightarrow C_J,
 \qquad k_p(z_i)=Z_{p(i)}.
\]
For nested ordered surjections, the maps compose by concatenating
the orders inside their fibres.
\end{proposition}
\begin{proof}
Localization is permitted by the preceding pairwise test.
The injection $e_p$ identifies each retained simplex with the
corresponding target simplex, including all its faces.
The component maps therefore preserve the defining compatibility
equations of $C_I$.
They preserve multiplication, the unit, and the polynomial-form
differential term by term.

Let $q:J\to K$ be another ordered surjection.
The order of $(qp)^{-1}(k)$ first lists the $J$-blocks in the
chosen order inside $q^{-1}(k)$, then their $I$-labels.
Thus $e_{qp}=e_p e_q$.
Both composite coefficient maps select the same component and
make the substitution $z_i\mapsto W_{q(p(i))}$.
Their simplex variables also have the same images.
This proves the identity $k_{qp}=k_qk_p$.
\end{proof}

The closed comparison form
$\omega_{ii'}=dq_{ii'}/(z_i-z_{i'})$ has a particularly simple image:
\begin{equation}\label{eq:moi-omega-image}
 k_p(\omega_{ii'})=
 \begin{cases}
 0,&p(i)=p(i'),\\
 \omega_{p(i),p(i')},&p(i)\ne p(i').
 \end{cases}
\end{equation}
In the first case, the selected simplex face makes $q_{ii'}$
constant, so its differential is zero before coordinates are
specialized.
In the second case its order indicator is the order indicator
of the two blocks, with the corresponding denominator.

For $p(1)=p(2)=a$, $p(3)=b$, and $1<2$ in the first block,
the retained charts are $123$ and $312$.
Thus $\omega_{12}$ vanishes, while both $\omega_{13}$ and
$\omega_{23}$ become $\omega_{ab}$.
Following this by $a<b$ selects $123$.
Merging $2<3$ first and then placing $1$ before that block
selects the same chart.

\section{The Koszul complex of a diagonal}

A collision sets coordinate differences to zero.
For each independent difference, adjoin a degree-minus-one exterior
generator whose differential is that difference.
The resulting Koszul complex computes derived restriction, as in
Chapter~\ref{chap:ordered-collisions-residues}.

Let $r_p(j)$ be the first label in the chosen order of $p^{-1}(j)$.
Introduce
\[
 x_i=z_i-z_{r_p(p(i))}\quad(i\ne r_p(p(i))).
\]
Using $Z_j=z_{r_p(j)}$ gives a polynomial identification
\[
 R_I=R_J[x_i\mid i\ne r_p(p(i))].
\]
The collision sets every $x_i$ to zero.

Form the graded-commutative algebra
\[
 K_p=R_I\otimes\Lambda(\epsilon_i\mid i\ne r_p(p(i))),
 \qquad |\epsilon_i|=-1,\qquad d\epsilon_i=x_i.
\]
The exterior generators anticommute and square to zero.
The differential fixes $R_I$ and obeys the graded Leibniz rule.
Its square vanishes on the generators, hence on the algebra.

For one normal coordinate, this is
$R_J[x]\epsilon\xrightarrow{x}R_J[x]$.
Multiplication by $x$ is injective, and its cokernel is $R_J$.
An explicit contraction over $R_J$ sends
\[
 f(x)\longmapsto
 \epsilon\,\frac{f(x)-f(0)}{x},
 \qquad \epsilon g(x)\longmapsto0.
\]
The quotient is polynomial.
On degree zero, $dh+hd$ is $f\mapsto f-f(0)$.
On degree minus one it is the identity.
Tensoring these one-variable contractions, with their tensor
signs, proves that $K_p\to R_J$ is a resolution.
It is bounded and free as a complex of $R_I$-modules.

For a complex $M$ of $R_I$-modules, the complex
$K_p\otimes_{R_I}M$ is its derived restriction to this diagonal.
Multiplication by $x_i$ has the explicit homotopy
\[
 x_i m=d(\epsilon_i m)+\epsilon_i\,dm.
\]
For a cycle $m$, this gives $x_i m=d(\epsilon_i m)$.
Exterior products of the $\epsilon_i$ supply the compatibility
between these homotopies.
This construction preserves quasi-isomorphisms.
Indeed, filter its tensor product with an acyclic complex by the
finite exterior degree. Each quotient is a finite sum of shifted
copies of that acyclic complex. Induction on this finite
filtration proves acyclicity.

Apply the construction to the coefficients:
\[
 B_p=K_p\otimes_{R_I}C_I.
\]
It is a commutative differential graded algebra over $R_J$,
using the representative-coordinate embedding just specified.
The previous map extends to
\begin{equation}\label{eq:moi-derived-coefficient}
 \kappa_p:B_p\longrightarrow C_J,\qquad
 \epsilon_i\longmapsto0,\qquad c\longmapsto k_p(c).
\end{equation}
This is a chain map because $k_p(x_i)=0$.
Thus the coefficient restriction exists on an explicit complex
computing the derived restriction. No flatness of $C_I$ has been
assumed.

The nested compatibility also has an explicit expression.
For $q:J\to K$, write
\[
 y_j=Z_j-Z_{r_q(q(j))},\qquad d\eta_j=y_j,
\]
with $y_j=\eta_j=0$ on representative labels.
For $p$ use $x_i=\epsilon_i=0$ on representatives.
A direct collision from $I$ to $K$ has normal coordinate
\[
 w_i=z_i-z_{r_p(r_q(q(p(i))))}=x_i+y_{p(i)}.
\]
Its exterior generator $\zeta_i$ is sent to
\begin{equation}\label{eq:moi-koszul-composition}
 \zeta_i\longmapsto\epsilon_i+\eta_{p(i)}.
\end{equation}
There are $|I|-|K|$ generators on each side.
The inverse sends
\[
 \eta_j\longmapsto\zeta_{r_p(j)},\qquad
 \epsilon_i\longmapsto\zeta_i-\zeta_{r_p(p(i))},
\]
where the final representative has $\zeta=0$.
These formulas are inverse linear changes of exterior generators.
They commute with the differential by the formula for $w_i$.
Consequently the complex for successive restrictions is
isomorphic to the complex for the direct restriction.
Every coefficient collision sends these exterior generators to zero.
The equality of their remaining maps is exactly $k_{qp}=k_qk_p$.
For further nested collisions, each generator becomes the sum along
its successive blocks, independent of parenthesization.

\section{Translation and composite states}

Translation of an insertion is differentiation in its position.
On a localized coefficient ring, define
\[
 \partial_{z_i}(z_i-z_{i'})^{-1}
 =-(z_i-z_{i'})^{-2}
\]
and extend by the usual product rule.
The derivations kill simplex variables.
They preserve face compatibility and commute with the coefficient
differential, hence act on $C_I$.

A connection on a module over $R_I$ means operators $\nabla_i$
satisfying
\[
 \nabla_i(fm)=(\partial_{z_i}f)m+f\nabla_i m.
\]
It is flat when the operators commute.
It is compatible with a complex when each commutes with its
differential.
Equivalently these operators make the module a module over the
algebra generated by $z_i,\partial_{z_i}$ with relations
\[
 [z_i,z_{i'}]=[\partial_{z_i},\partial_{z_{i'}}]=0,\qquad
 [\partial_{z_i},z_{i'}]=\delta_{ii'}.
\]
This algebra is the algebra of polynomial differential operators
on the coordinate space. These are left differential modules.

The tangent derivative along a merged coordinate is
\[
 \nabla_{Z_j}=\sum_{p(i)=j}\partial_{z_i}.
\]
It annihilates every $x_i$.
On $B_p$, let it also annihilate $\epsilon_i$.
It then commutes with $d\epsilon_i=x_i$.
The chain rule gives
\begin{equation}\label{eq:moi-connection}
 \partial_{Z_j}\kappa_p(c)
 =\kappa_p\left(\sum_{p(i)=j}\partial_{z_i}c\right).
\end{equation}
Thus collision respects translation along the merged point.
For nested collisions, summing tangent derivatives over blocks
gives the derivative for the composite block.

Now introduce an actual composite state for every finite word.
Let $V_I$ have basis $u_i$, all of degree zero and weight one.
Its tensor algebra is the vector space of finite linear
combinations of words, with concatenation as multiplication.
Set
\[
 \mathcal A_I=C_I\otimes_{\mathbb C}T(V_I).
\]
The differential acts on coefficients. The states are horizontal:
$\nabla_i u_{i'}=0$.
The augmentation kills every nonempty word.
The algebra is associative because concatenation is associative.

After extending coefficients by \eqref{eq:moi-derived-coefficient},
define
\[
 u_{i_1}\cdots u_{i_m}
 \longmapsto u_{p(i_1)}\cdots u_{p(i_m)}.
\]
It preserves products, differential, weight, augmentation, and
the connection. In particular, $u_1u_2$ becomes the nonzero state
$u_a^2$ when $1,2$ merge to $a$.
The target retains a composite word at one position.

This free-word algebra differs from the finite algebra $A_I$
of Chapter~\ref{chap:ordered-coefficient-bar}.
There the coefficient $\omega_{ij}$ entered the state product.
Here products of states are horizontal, while collision classes
remain in the coefficient algebra. Both constructions have
explicit bar differentials. An identification between them would
require an additional map preserving the structures in question.

\section{Bar words follow the merged states}

Put $Q_I=B_{C_I}\mathcal A_I$.
Its elements are words $[a_1|\cdots|a_\ell]$ in nonempty states,
with $|sa|=|a|-1$.
The differential has internal part $b_1(sa)=-s(da)$ and
multiplication part
\[
 b_2(sa\otimes sb)=(-1)^{|a|}s(ab).
\]
These are the conventions already derived in
Chapter~\ref{chap:ordered-coefficient-bar}.
Disjoint adjacent contractions cancel by their odd degree.
Overlapping contractions give the signed associator, and mixed
terms give the coefficient Leibniz identity.
Those same calculations prove $b^2=0$ for $\mathcal A_I$.

For a $C_I$-module $M$, define its coefficient extension by
\[
 E_p(M)=C_J\otimes_{B_p}^{\mathbb L}
                  (K_p\otimes_{R_I}M).
\]
The superscript means that a free resolution is used before
tensoring when necessary.
For the bar modules at a fixed weight, no unspecified resolution
is needed. A word of weight $w$ uses at most $w$ letters.
There are finitely many label words and cuts at that weight.
Filtering by bar length makes $Q_I(w)$ a finite semifree
$C_I$-module: the differential of each state-word basis vector
uses shorter bar words.
The same property holds after tensoring with $K_p$ over $R_I$.
Thus ordinary scalar extension computes $E_p(Q_I(w))$.

The collision map on the relative bar is
\[
 \gamma_p:E_p(Q_I)\longrightarrow Q_J,\qquad
 [a_1|\cdots|a_\ell]\longmapsto
 [\alpha_p(a_1)|\cdots|\alpha_p(a_\ell)].
\]
Here $\alpha_p$ denotes the coefficient and state map just defined.
It preserves each product and each degree.
Substitution in $b_1,b_2$ shows that $\gamma_p$ is a chain map.
It also preserves every word cut.
All these maps respect the connection by
\eqref{eq:moi-connection}.
The explicit changes of exterior generators in
\eqref{eq:moi-koszul-composition} give the same compatibility
for nested collisions.

For the three-label example,
\[
 b[u_1|u_2|u_3]=[u_1u_2|u_3]-[u_1|u_2u_3].
\]
After merging $1,2$ this becomes
\[
 [u_a^2|u_b]-[u_a|u_au_b].
\]
Its next boundary is zero by associativity.
Multiplying the original word by $\omega_{13}$ transports that
coefficient to $\omega_{ab}$ with the same tensor signs.
Thus state merger and coefficient restriction occur in one
nonzero bar calculation.

The resolution $\mathcal A_I\otimes_{C_I}Q_I$ has leading face
\[
 a_0\otimes[a_1|\cdots|a_\ell]\longmapsto
 -(-1)^{|a_0|}a_0a_1\otimes[a_2|\cdots|a_\ell].
\]
The insertion $-1\otimes[\overline a_0|a_1|\cdots|a_\ell]$
contracts its augmentation.
The cancellation and semifree lifting argument in the preceding
bar chapter depend only on associativity, augmentation, and a
free positive-weight state basis. They apply here at each weight
and to the full algebraic resolution.
Consequently its coefficient dual computes
$R\operatorname{Hom}_{\mathcal A_I}(C_I,C_I)$.
The product is convolution by cuts, equivalently composition of
the lifted resolution endomorphisms. This is a relative algebraic
Hom, not the space of horizontal maps.

\section{Why dual collision changes the coefficients}

Write $Q_I^\vee(w)=\operatorname{Hom}_{C_I}(Q_I(w),C_I)$.
Precomposition with $\gamma_p$ defines
\begin{equation}\label{eq:moi-dual-collision}
 d_p^\vee:Q_J^\vee(w)\longrightarrow E_p(Q_I^\vee(w)).
\end{equation}
To justify this target, consider one shifted free $C_I$-module.
Evaluation identifies its dual after extension with the extension
of its dual. The identification persists under finite direct
sums, shifts, and cones.
The finite bar-length filtration therefore gives
\[
 \operatorname{RHom}_{C_J}(E_p(Q_I(w)),C_J)
 \simeq E_p(Q_I^\vee(w)).
\]
Differentiating a coefficient-linear map requires subtracting the
derivative of its input. For complexes $M,N$ over $C_J$ with the
flat tangent connections just constructed, define
\[
 \nabla_j^{\operatorname{Hom}}F
      =\nabla_j^N\circ F-F\circ\nabla_j^M,
 \qquad
 d_{\operatorname{Hom}}F=d_N\circ F-(-1)^{|F|}F\circ d_M.
\]
Here $F$ is homogeneous and the connections have degree zero.
For homogeneous $c\in C_J$, the two terms differentiating $c$
cancel. Thus
\[
 (\nabla_j^{\operatorname{Hom}}F)(cm)
       =(-1)^{|F||c|}c(\nabla_j^{\operatorname{Hom}}F)(m),
\]
so this is again a graded coefficient-linear map.
Expanding two commutators gives
\[
 [\nabla_i^{\operatorname{Hom}},\nabla_j^{\operatorname{Hom}}]F
   =[\nabla_i^N,\nabla_j^N]F-F[\nabla_i^M,\nabla_j^M]=0.
\]
Since each connection commutes with its complex differential,
expansion of the displayed Hom differential also gives
$d_{\operatorname{Hom}}\nabla_j^{\operatorname{Hom}}
 =\nabla_j^{\operatorname{Hom}}d_{\operatorname{Hom}}$.

Under the duality identification above, the map in
\eqref{eq:moi-dual-collision} is $F\mapsto F\gamma_p$.
The chain and horizontal identities for $\gamma_p$ imply
\[
 d_{\operatorname{Hom}}(F\gamma_p)
     =(d_{\operatorname{Hom}}F)\gamma_p,
 \qquad
 \nabla_j^{\operatorname{Hom}}(F\gamma_p)
     =(\nabla_j^{\operatorname{Hom}}F)\gamma_p.
\]
Thus dual collision preserves both the differential and connection.

Complete by weight in each cohomological degree, retaining a
direct sum over cohomological degrees.
The dual map then has target
\[
 \prod_{w\ge0} E_p(Q_I^\vee(w)).
\]
Coefficients are extended at each finite weight before taking
this product. No interchange of an unrestricted product with
derived restriction is asserted.
Convolution uses finitely many splittings of each output weight.
The coalgebra identity for $\gamma_p$ therefore makes the
completed dual map unital and multiplicative.
For nested collisions the composition is
\[
 d_{qp}^\vee=E_q(d_p^\vee)\circ d_q^\vee,
\]
under the coefficient-extension identification supplied by the
isomorphism between the Koszul complexes for nested restrictions.

The simplest coefficient fibre forces this extension.
Take $B=\mathbb C^2$, $C=\mathbb C$, and let $B$ act on $C$
by the first projection.
If $f:C\to B$ is $B$-linear, then
\[
 (0,1)f(1)=f((0,1)\cdot1)=0.
\]
Thus $f(1)$ lies in the first summand and cannot equal $(1,1)$.
There is no unital reverse $B$-linear map from $C$ to $B$.
After extension, $C\otimes_B B=C$, and the weight-zero
dual map is the identity of $C$.

The completed cobar counit has the corresponding covariance.
It sends $t(sa)$ to $a$ and longer bar generators to zero.
Its internal and cut terms cancel on a two-letter generator.
The finite gap contraction proved in
Theorem~\ref{thm:ocb-reconstruction} applies at every weight:
each original state has positive weight, and only finitely many
letters and cuts can occur.
It gives
\[
 \widehat\Omega_{C_I}Q_I\simeq\widehat{\mathcal A}_I.
\]
The target is now completed as well. For example,
$\sum_{m\ge0}u_i^m$ is a degree-zero element of this target.
Both the differential and counit commute with collision because
collision preserves products and cuts.
All completed statements are obtained weight by weight.

\section{What order retains at the diagonal}

Fix the merged coordinate in the binary case and write $x$ for
the separation. The coefficient complex is
\[
 [R^2\xrightarrow{(a,b)\mapsto b-a}R[x^{-1}]],
 \qquad R=\mathbb C[x].
\]
Its comparison with the compatible polynomial-form model is explicit.
A degree-zero path has coefficients in $R[x^{-1}][t]$ and endpoints
in $R$; a degree-one path is $g(t)dt$ with localized coefficients.
Send the path to its two endpoints and send $g(t)dt$ to
$\int_0^1g(t)dt$.
The fundamental theorem for polynomial integration makes this a
map of complexes to the displayed two-term complex.
It is surjective: interpolate the two endpoints linearly, and
use a constant one-form for any prescribed integral.
Its kernel contracts by $g(t)dt\mapsto\int_0^t g(s)ds$.
The zero-integral condition gives the second zero endpoint.
Thus it is a quasi-isomorphism of complexes over $R$.
No multiplicativity of polynomial integration is needed.
Its terms are flat over $R$.
Its derived fibre at $x=0$ is therefore $\mathbb C^2$ in degree
zero, with its two coordinate projections.
Endpoint evaluation is itself multiplicative, so this identifies
the product on that fibre with coordinatewise multiplication.
The coefficient algebra for one merged label has fibre
$\mathbb C$. Neither projection is an equivalence.

This calculation separates the ordered family from descent that
forgets the number of repeated labels.
Such descent would require a family on finite powers whose
derived restrictions along repeated-coordinate maps are
equivalences to the merged family.
This is the Cartesian diagonal condition used for the usual
space of finite subsets, called the Ran space.
The binary ranks already show that the present ordered family
does not satisfy that condition.
The constructed collision maps retain the chosen order precisely
because they are noninvertible restrictions.

There is also an obstruction to declaring every earlier state
model horizontal.
For the finite coefficient-weighted multiplication
$a_i a_j=\omega_{ij}b_{ij}$, suppose
$\nabla a_i=\nabla a_j=\nabla b_{ij}=0$.
Then the connection Leibniz rule would require
\[
 0=\partial_{z_i}(\omega_{ij}b_{ij})
   =-\frac{dq_{ij}}{(z_i-z_j)^2}b_{ij}.
\]
The right side is nonzero on the adjacent-exchange simplex.
The coefficient-weighted algebra therefore does not acquire
this horizontal connection.
A differential-module realization would need additional state
connection data or a different explicitly compared construction.

We have obtained coherent coefficient restrictions, composite
states at merged points, and their bar and dual maps on finite
powers. Forgetting order, realizing coefficient-weighted states,
and comparing with singular operator expansions each require
further maps. Their domains and the first obstructing identities
are now explicit.
